\documentclass{article}
\usepackage[utf8]{inputenc}
\title{Olá, mundo!}
\author{Bryan Samuel da Silva Cherobini}
\date{16 de Abril de 2026}

\begin{document}

\maketitle

\section{Introdução}


Olá, mundo!
Este documento apresenta a resolução do exercício 1 com o objetivo de revisar a estrutura básica da linguagem LaTeX e a introdução e familiarização do autor com a linguagem.

\section{Desenvolvimento}

Este documento apresenta a resolução do exercício 1 com o objetivo de revisar a estrutura básica da linguagem LaTeX e promover a minha introdução e familiarização com a ferramenta. Por ser o meu primeiro arquivo de autoria própria, ele funciona como um ponto de partida para entender como os comandos funcionam na prática, transformando o código LaTeX em um artigo organizado e estruturado.

O foco principal desse artigo é a adaptação ao novo formato de escrita e linguagem, trocando os editores comuns por uma linguagem de marcação. Além de servir como uma introdução técnica, este trabalho me ajuda a ganhar confiança na criação de documentos originais e garante que eu consiga dominar as bases necessárias para os próximos desafios acadêmicos que utilizarem o LaTeX.

\section{Conclusão}

Pode-se concluir que este artigo serviu como uma base essencial de conhecimento e introdução à escrita em LaTeX, cumprindo seu papel de ambientação técnica. Ao finalizar este documento, consolidei o aprendizado sobre a estruturação de títulos, autores, seções e subseções, garantindo uma organização visual clara. Além disso, mesmo que recursos como tabelas não tenham sido utilizados na prática neste primeiro momento, a revisão teórica sobre o tema foi fundamental para a minha adaptação e familiarização completa com a linguagem, preparando o caminho para a criação de trabalhos originais mais complexos.

\end{document}